\begin{frame}{R-CNN (2014): ``후보 $\to$ 특징 $\to$ 판정'' 3단계}
  \begin{columns}[T]
    \begin{column}{0.56\textwidth}
      \centering
      \includegraphics[width=\linewidth]{ch10_3_rcnn_pipeline.png}
      \\[0.3em]{\footnotesize R-CNN 파이프라인 ---
      (1) 셀렉티브 서치로 $\sim$2,000개 후보 박스,
      (2) 각 후보를 224$\times$224로 줄여 CNN을 통과시켜
      4,096차원 특징, (3) 클래스별 SVM + 박스 회귀로 판정,
      (4) IoU 기반 NMS로 중복 박스 제거.}
    \end{column}
    \begin{column}{0.44\textwidth}
      \small
      \begin{enumerate}
        \item \kb{후보 지역}: \textbf{셀렉티브 서치}
          (2012, \textbf{딥러닝이 아닌 기술})가 색·질감이
          비슷한 인접 영역을 합치며 후보 박스
          \textbf{$\sim$2,000개}. 슬라이딩 윈도우(수백만 개)
          대비 ``미리 좁히기''가 통하는 이유: 사물은 실제로
          그렇게 많지 않고 작지도 않은 경우가 대부분
        \item \kb{특징 추출}: 2,000개 후보를 각각
          224$\times$224로 비율만 맞춰 리사이즈하고,
          \textbf{분류용 CNN의 마지막 분류층을 제외한
          부분}을 그대로 통과 --- 각 후보마다 4,096차원
          특징 벡터 1개. \textbf{``분류로 배운 특징
          추출기''를 새 문제에 재사용하는} 순간
        \item \kb{판정}: (a) \textbf{클래스별 SVM 하나씩}
          (Week 5의 SVM --- 20개 클래스면 SVM 20개)로
          ``무엇이'', (b) \textbf{회귀 모델}로 ``정확한
          사물 경계로 미세 조정''. 2,000 후보 $\times$ SVM 20개
          $= 40{,}000$회 판정 --- ``한 특징 벡터는 1번만
          계산해도 4만 번 재사용''되는 또 하나의 파라미터
          재사용
      \end{enumerate}
    \end{column}
  \end{columns}
\end{frame}
